\chapter{System overview}
\label{ch:soc}

\section{Delivered scope}

The delivery is three Verilog modules and the verification environment around
them:

\begin{itemize}[nosep]
  \item \file{hdl/cpu\_top.v} --- RV32I core, three pipeline stages, two AXI4-Lite
        master ports.
  \item \file{hdl/pic.v} --- interrupt controller, 16 hardware lines plus 16
        software channels, AXI4-Lite slave.
  \item \file{hdl/mtimer.v} --- 64-bit machine timer, AXI4-Lite slave.
\end{itemize}

The CPU was specified by a feature brief. The PIC implements a second brief that
no block in the wider project owned. The machine timer had no brief at all: it
was added because without a timer the system has no deadline mechanism for a
stalled DMA transfer and no way to wake the core from \sig{WFI}.

The remaining SoC blocks --- multi-channel DMA, dual-port SRAM, interconnect ---
are outside this delivery. They appear below only as the environment the CPU and
PIC were designed to meet.

\section{Block structure}

\begin{figure}[H]
\centering
\begin{tikzpicture}[x=1mm,y=1mm]
  % --- blocks -------------------------------------------------------------
  \node[blkf,minimum width=26,minimum height=11] (imem) at (26,68)
        {instruction memory};
  \node[blk,minimum width=36,minimum height=15]  (cpu)  at (26,50)
        {\textbf{cpu\_top}\\[1pt]\scriptsize RV32I, 3 stages\\\scriptsize 2 AXI4-Lite masters};
  \node[blkf,minimum width=28,minimum height=12] (per)  at (86,50)
        {DMA, DP-SRAM\\\scriptsize (not delivered)};
  \node[blkf,minimum width=30,minimum height=11] (fab)  at (26,30)
        {interconnect\\\scriptsize address decode};

  \node[blk,minimum width=26,minimum height=13]  (pic)  at (2,6)
        {\textbf{pic}\\\scriptsize 0x3000\_0000};
  \node[blk,minimum width=22,minimum height=13]  (tmr)  at (38,6)
        {\textbf{mtimer}\\\scriptsize 0x3001\_0000};
  \node[blk,minimum width=26,minimum height=13]  (dmem) at (74,6)
        {data memory\\\scriptsize / DP-SRAM port A};

  % --- instruction and data buses -----------------------------------------
  \draw[flow,<->] (imem.south) -- node[lbl,right=1pt]{ibus (AR/R)} (cpu.north);
  \draw[flow,<->] (cpu.south)  -- node[lbl,right=1pt]{dbus} (fab.north);

  % peripheral bus rail: one horizontal run, one drop per slave
  \draw[flow,-] (fab.south) -- (26,20);
  \draw[thick,draw=accent] (2,20) -- (74,20);
  \foreach \x in {2,38,74} \draw[flow] (\x,20) -- (\x,12.6);

  % --- interrupt sources into the PIC (routed below every block) ----------
  \draw[flow] ($(per.south)+(9,0)$) -- (95,-13) -- (6,-13) -- (6,-0.6);
  \node[lbl] at (46,-13) {irq lines from the peripherals};
  \draw[flow] (tmr.south) -- (38,-6.5) -- (-4,-6.5) -- (-4,-0.6);
  \node[lbl] at (20,-6.5) {irq $\to$ source 7};

  % --- CPU <-> PIC handshake, in a channel to the left of every block -----
  \draw[flow] (-11,6) -- (-20,6) -- (-20,50) -- (8,50);
  \node[lbl,rotate=90,align=center] at (-23,30) {cpu\_irq,\\cpu\_irq\_vec};
  \draw[flow] (8,45) -- (-15,45) -- (-15,2) -- (-11,2);
  \node[lbl,rotate=90] at (-17.5,26) {ack / eoi};
\end{tikzpicture}
\caption{Delivered blocks (solid outline) and their environment. Addresses are the
values used by the verification environment, not a fixed system memory map.}
\label{fig:soc}
\end{figure}

Everything runs in one clock domain with a single asynchronous, active-low reset.
The CPU's two master ports are independent: an instruction fetch and a data
access proceed at the same time, each with at most one transaction outstanding.

\section{AXI4-Lite as used in this system}
\label{sec:axiprimer}

Every bus in the delivery is AXI4-Lite. This section states the parts of the
protocol the rest of the document relies on, and what this system does and does
not use.

\subsection{Channels}

A transfer is carried on one of five independent channels. Each has its own
handshake and none of them blocks another.

\begin{table}[H]
\centering
\caption{AXI4-Lite channels and their payloads}
\label{tab:axichan}
\begin{tabular}{@{}llL{0.44\linewidth}@{}}
\toprule
\textbf{Channel} & \textbf{Direction} & \textbf{Carries} \\
\midrule
AW & master $\to$ slave & Write address (\sig{AWADDR}) and protection type (\sig{AWPROT}). \\
W  & master $\to$ slave & Write data (\sig{WDATA}) and byte strobes (\sig{WSTRB}). \\
B  & slave $\to$ master & Write response (\sig{BRESP}). \\
AR & master $\to$ slave & Read address (\sig{ARADDR}), protection type (\sig{ARPROT}). \\
R  & slave $\to$ master & Read data (\sig{RDATA}) and read response (\sig{RRESP}). \\
\bottomrule
\end{tabular}
\end{table}

A read uses AR then R. A write uses AW and W --- which are separate channels and
may be accepted in either order --- and is completed by B. There is no single
``write done'' signal other than the B response, which is why the write path in
this design tracks AW and W independently.

\subsection{The handshake}

Every channel works the same way. The source raises \sig{VALID} when it has
something to transfer; the destination raises \sig{READY} when it can take it.
\textbf{The transfer happens on the rising clock edge where both are high.} Either
signal may lead the other.

\begin{figure}[H]
\centering
\begin{tikztimingtable}[timing/slope=0,timing/coldist=3pt,xscale=1.5]
  clk    & 12{C} \\
  VALID  & 2L 2H 2L 4H 2L \\
  READY  & 2L 2H 4L 2H 2L \\
  PAYLOAD& 2U 2D{A} 4D{B} 2D{B} 2U \\
  \extracode
  \begin{pgfonlayer}{background}
    \vertlines[help lines]{2,4,6,8,10}
  \end{pgfonlayer}
\end{tikztimingtable}
\caption{Two transfers on one channel. The first (payload~A) completes
immediately because both sides are ready in the same cycle. The second
(payload~B) is held by the source for two extra cycles while the destination is
not ready; the payload does not change during the wait.}
\label{fig:axihs}
\end{figure}

Two rules follow, and both are checked on every bus in this system by the
protocol monitors and the assertion layer:

\begin{enumerate}[nosep]
  \item \textbf{\sig{VALID} must not wait for \sig{READY}.} A source that only
        asserts \sig{VALID} after seeing \sig{READY} deadlocks against a
        destination that only asserts \sig{READY} after seeing \sig{VALID}.
  \item \textbf{Once \sig{VALID} is high, the payload must not change} until the
        handshake completes. The destination is entitled to sample it in any
        cycle where both are high.
\end{enumerate}

A third rule matters at reset: \sig{VALID} must be low while reset is asserted.
A violation of exactly this rule was the first real defect the assertion layer
found in this project (Chapter~\ref{ch:verif}).

\subsection{Responses}

\begin{table}[H]
\centering
\caption{Response codes on \sig{RRESP} and \sig{BRESP}}
\label{tab:axiresp}
\begin{tabular}{@{}lllL{0.40\linewidth}@{}}
\toprule
\textbf{Code} & \textbf{Name} & \textbf{Source} & \textbf{Meaning here} \\
\midrule
\texttt{00} & OKAY   & slave        & Normal completion. \\
\texttt{01} & EXOKAY & ---          & Does not exist in AXI4-Lite. The monitors treat it as an error if it ever appears. \\
\texttt{10} & SLVERR & slave        & A real slave refused the access: a read-only register written, an unmapped offset inside a peripheral. \\
\texttt{11} & DECERR & interconnect & No slave is mapped at that address. \\
\bottomrule
\end{tabular}
\end{table}

The CPU turns both error codes into precise traps rather than ignoring them ---
cause 1 on an instruction fetch, cause 5 or 7 on a load or store
(Chapter~\ref{ch:csr}).

\subsection{What this system does not use}

AXI4-Lite is already the reduced profile: no bursts, no transaction identifiers,
no out-of-order completion, and all transfers are the full data width with byte
strobes selecting lanes. On top of that, this design constrains itself further:

\begin{itemize}[nosep]
  \item \textbf{At most one transaction outstanding per port.} A master issues an
        address and waits for the response before issuing another.
  \item \textbf{Never a read and a write in flight together} on the data bus.
  \item \textbf{No write interleaving} and no need for reordering logic anywhere.
\end{itemize}

These are design choices, not protocol requirements, and they are what make the
masters small. The per-channel handshake tracking described in
Chapter~\ref{ch:lsu} was written so that lifting the one-outstanding restriction,
or moving to AXI4-Full, extends the existing structure rather than replacing it.

\section{Port naming convention}
\label{sec:portnaming}

Every port of every module in this delivery carries a direction suffix:
\texttt{\_i} on an input, \texttt{\_o} on an output. The rule has no exceptions,
including the clock and reset (\sig{clk\_i}, \sig{rst\_n\_i}) and the AXI4-Lite
channels (\sig{s\_axi\_awaddr\_i}, \sig{s\_axi\_bresp\_o}).

The point is that direction becomes readable at the point of use rather than
only at the declaration. In a port map several hundred lines away from the
module header, \sig{s\_axi\_awready\_o} states which side drives the wire without
opening another file, and a connection that has the direction backwards is
visible as a mismatched pair rather than as an elaboration error.

\begin{table}[H]
\centering
\caption{The convention, and what it looks like on each kind of port}
\label{tab:portnaming}
\begin{tabular}{@{}L{0.22\linewidth}L{0.28\linewidth}L{0.36\linewidth}@{}}
\toprule
\textbf{Kind of port} & \textbf{Example} & \textbf{Note} \\
\midrule
Clock and reset & \sig{clk\_i}, \sig{rst\_n\_i} & The reset stays active low; the
  suffix says only that the module receives it. \\
Scalar or bus signal & \sig{cpu\_irq\_i}, \sig{cpu\_irq\_ack\_o} & The same
  logical signal is \sig{cpu\_irq\_o} on the PIC and \sig{cpu\_irq\_i} on the
  CPU, which is exactly the information the suffix is there to carry. \\
AXI4-Lite master port & \sig{dbus\_axi\_awaddr\_o}, \sig{dbus\_axi\_awready\_i} &
  Address, data and \sig{VALID} out; \sig{READY} and the response in. \\
AXI4-Lite slave port & \sig{s\_axi\_awaddr\_i}, \sig{s\_axi\_awready\_o} & The
  mirror of the master port. \\
\bottomrule
\end{tabular}
\end{table}

Internal signals are \emph{not} suffixed. A wire inside a module has no
direction, so a suffix on one would be noise, and the absence of a suffix is
itself a useful signal to the reader: an identifier without one is local.

\begin{deviation}
\textbf{Consequence for AXI4-Lite tooling.} The AXI channel names in this
delivery are the protocol names plus the direction suffix, so they no longer
match the exact spelling ARM IHI0022E uses. Vendor flows that infer a bus
interface from port names alone (IP-XACT generation, Vivado's automatic AXI
inference) will not recognise them without a name mapping. The trade was made
deliberately, in favour of uniform readability across the whole delivery; the
channel structure, widths and behaviour are unchanged, so a mapping is a
rename table and nothing more.
\end{deviation}

\section{Interrupt path}

The interrupt path is the only place where all three delivered blocks meet, and
it is worth stating once because the rest of the document refers back to it.

Peripherals raise level-sensitive lines into the PIC. The PIC resolves priority
itself and presents the CPU with a \emph{single} request plus the identifier of
the winning source. The CPU accepts it only at an instruction boundary and only
if both its own \reg{mie} bit and the global \reg{mstatus.MIE} allow it. Entering
the handler produces a one-cycle \sig{cpu\_irq\_ack} pulse --- the PIC's
\emph{claim} --- and the interrupt-returning \sig{MRET} produces a one-cycle
\sig{cpu\_irq\_eoi} pulse --- the \emph{end of interrupt}.

\begin{table}[H]
\centering
\caption{CPU\,$\leftrightarrow$\,PIC interface (\file{hdl/cpu\_top.v}, \file{hdl/pic.v})}
\label{tab:cpupic}
\begin{tabular}{@{}L{0.32\linewidth}lL{0.38\linewidth}@{}}
\toprule
\textbf{Signal} & \textbf{Direction} & \textbf{Function} \\
\midrule
\sig{cpu\_irq\_o} $\to$ \sig{cpu\_irq\_i} & PIC $\to$ CPU & Level request. High while the PIC is offering a source. \\
\sig{cpu\_irq\_vec\_o[3:0]} $\to$ \sig{cpu\_irq\_vec\_i} & PIC $\to$ CPU & Identifier of the offered source; becomes trap cause $16+v$. \\
\sig{cpu\_irq\_ack\_o} $\to$ \sig{cpu\_irq\_ack\_i} & CPU $\to$ PIC & One-cycle claim pulse, asserted when the CPU commits to the handler. \\
\sig{cpu\_irq\_eoi\_o} $\to$ \sig{cpu\_irq\_eoi\_i} & CPU $\to$ PIC & One-cycle pulse on an interrupt-returning \sig{MRET}; pops one nesting level. \\
\sig{cpu\_in\_trap\_o} & CPU $\to$ observer & High from trap entry until \sig{MRET}. Not used by the PIC; kept for observability and assertions. \\
\bottomrule
\end{tabular}
\end{table}

The design carries a consequence worth flagging at system level: a source that is
enabled in the PIC but masked in the CPU's \reg{mie} stays the PIC's offer
indefinitely, because nothing ever claims it, and starves every less urgent
source behind it. Deadline escalation (Section~\ref{sec:deadline}) does not help
here --- the problem is that nobody claims, not that the source is insufficiently
urgent. The integration rule is therefore: do not route a source into the PIC
unless some handler will service it.

\section{Address map used in verification}

The system memory map is not fixed. The reset vector is the \sig{RESET\_PC}
parameter and the peripheral bases are decoder parameters in the testbench,
collected in \file{debug/sim/soc\_map.vh}.

\begin{table}[H]
\centering
\caption{Address map of the verification environment}
\label{tab:memmap}
\begin{tabular}{@{}llL{0.50\linewidth}@{}}
\toprule
\textbf{Region} & \textbf{Base} & \textbf{Notes} \\
\midrule
Instruction memory & \texttt{0x0000\_0000} & Model, 1024 words, initialised from the assembled program. \\
Data memory        & \texttt{0x0000\_2000} & Model, 1024 words; default leg of the decoder. \\
PIC                & \texttt{0x3000\_0000} & 56 mapped words, see Appendix~\ref{app:regmap}. \\
Machine timer      & \texttt{0x3001\_0000} & 4 mapped words. \\
\bottomrule
\end{tabular}
\end{table}

Both peripheral slaves decode only \sig{addr[7:2]}. Inside an interconnect window
larger than 256\,bytes, higher addresses therefore alias onto the first registers
instead of producing a decode error. This is recorded as a limitation in
Chapter~\ref{ch:limits}.
